\documentclass[11pt]{article}
% ======================================================================
%  examstyle.tex — EPFL-style exam layout for the weekly Time Series exams
%  Shared by every Exam_Week_NN(.tex / _Solutions.tex). Compiles with pdflatex.
% ======================================================================
\usepackage[utf8]{inputenc}
\usepackage[T1]{fontenc}
\usepackage[english]{babel}
\usepackage[a4paper,margin=2.4cm]{geometry}
\usepackage{amsmath,amssymb,amsthm,mathtools}
\usepackage{bm}
\usepackage{enumitem}
\usepackage{array}

\setlength{\parindent}{0pt}
\setlength{\parskip}{0.5em}
\emergencystretch=3em
\allowdisplaybreaks[1]

% ---------- math macros (same names as the rest of the project) ----------
\newcommand{\E}{\mathbb{E}}
\DeclareMathOperator{\Var}{Var}
\DeclareMathOperator{\Cov}{Cov}
\DeclareMathOperator{\Corr}{Corr}
\DeclareMathOperator*{\argmin}{arg\,min}
\DeclareMathOperator{\MSE}{MSE}
\DeclareMathOperator{\trace}{tr}
\newcommand{\Reals}{\mathbb{R}}
\newcommand{\Ints}{\mathbb{Z}}
\newcommand{\Comp}{\mathbb{C}}
\newcommand{\Nat}{\mathbb{N}}
\newcommand{\Prob}{\mathbb{P}}
\newcommand{\Tr}{^{\mathsf{T}}}
\newcommand{\conj}[1]{\overline{#1}}
\newcommand{\abs}[1]{\left\lvert #1 \right\rvert}
\newcommand{\norm}[1]{\left\lVert #1 \right\rVert}
\newcommand{\ind}{\mathbf{1}}
\newcommand{\eps}{\varepsilon}
\newcommand{\diff}{\nabla}
\newcommand{\acvs}{\gamma}
\newcommand{\acf}{\rho}
\newcommand{\pacf}{\alpha}
\newcommand{\sdf}{\mathcal{S}}
\newcommand{\Bsh}{B}
\newcommand{\iid}{\overset{\text{iid}}{\sim}}

\setlist[enumerate]{leftmargin=2em,topsep=2pt,itemsep=4pt}
\setlist[itemize]{leftmargin=1.6em,topsep=2pt,itemsep=2pt}

\newcounter{exo}

% ---------- exam cover header :  \examheader{exam title}{duration} ----------
\newcommand{\examheader}[2]{%
  \thispagestyle{empty}
  \begin{center}
    {\Large\bfseries Time Series}\\[3pt]
    Spring Semester 2026, EPFL\\
    \'Ecole Polytechnique F\'ed\'erale de Lausanne\\
    Prof.\ Dr.\ Sofia C.\ Olhede\\[7pt]
    \hrule height 0.8pt
    \vspace{9pt}
    {\large\bfseries Practice Exam --- #1}\\[4pt]
    {\normalsize Duration: #2 \quad$\cdot$\quad No calculator \quad$\cdot$\quad Answer on paper}\\
    \vspace{8pt}
    \hrule height 0.8pt
  \end{center}
  \vspace{14pt}
  \begin{tabular}{@{}p{8.2cm}@{}p{8cm}@{}}
    Name:\enspace\hrulefill & Surname:\enspace\hrulefill
  \end{tabular}\\[14pt]
  SCIPER (Student Card Number):\enspace\hrulefill
  \vspace{16pt}
}

% ---------- points table (fixed: 4 problems) ----------
\newcommand{\pointstable}{%
  \begin{center}
  \renewcommand{\arraystretch}{1.5}
  \begin{tabular}{|p{5cm}|p{4cm}|}
    \hline
    \textbf{Exercise} & \textbf{Points}\\\hline
    1 & \\\hline
    2 & \\\hline
    3 & \\\hline
    4 & \\\hline
    \textbf{Total} & \\\hline
  \end{tabular}
  \end{center}
}

% ---------- problem heading (exam: one problem per page) ----------
\newcommand{\problem}[1]{%
  \clearpage
  \stepcounter{exo}%
  \begin{center}\textbf{\large Problem \theexo\ (#1 points)}\end{center}
  \vspace{4pt}
}
% ---------- answer space: fill the statement page + one blank answer page ----------
% Put \answerpage at the END of each problem so every exercise gets its own page(s).
\newcommand{\answerpage}{\par\vfill\newpage\null\newpage}

% ---------- solutions document ----------
\newcommand{\solheader}[1]{%
  \begin{center}
    {\Large\bfseries Time Series --- Solutions}\\[3pt]
    {\large Practice Exam --- #1}\\[2pt]
    EPFL, MATH-342 \quad$\cdot$\quad Prof.\ S.\ C.\ Olhede\\[5pt]
    \hrule height 0.8pt
  \end{center}
  \vspace{10pt}
}
\newcommand{\solproblem}[1]{%
  \bigskip
  \stepcounter{exo}%
  {\large\bfseries Problem \theexo\ (#1 points)}\par\nobreak\vspace{2pt}
}
% Rappel de l'énoncé avant la solution (dans les corrigés)
\newcommand{\statementstart}{\par\vspace{1pt}{\bfseries\itshape Statement.}\par\nobreak\begingroup\small\itshape}
\newcommand{\statementend}{\par\endgroup\vspace{5pt}\hrule\vspace{6pt}{\bfseries Solution.}\par\nobreak\vspace{2pt}}

% --- local macros not in examstyle ---
\newcommand{\F}{\mathcal{F}}
\begin{document}

\solheader{Week 12: Non-Linear Time Series --- ARCH \& GARCH}

% ---------------------------------------------------------------
\solproblem{5}
\textbf{(a)} Since $\sigma_t$ is $\F_{t-1}$-measurable and $\eta_t$ is independent of $\F_{t-1}$ with $\E[\eta_t]=0$,
\[
\E[r_t\mid\F_{t-1}]=\E[\sigma_t\eta_t\mid\F_{t-1}]=\sigma_t\,\E[\eta_t]=\sigma_t\cdot0=0 .
\]
By the law of iterated (tower) expectation,
\[
\E[r_t]=\E\bigl[\E[r_t\mid\F_{t-1}]\bigr]=\E[0]=0 .
\]
The conditional variance is, since $\E[r_t\mid\F_{t-1}]=0$ and $\Var(\eta_t)=1$,
\[
\Var(r_t\mid\F_{t-1})=\E[r_t^2\mid\F_{t-1}]=\sigma_t^2\,\E[\eta_t^2\mid\F_{t-1}]=\sigma_t^2
=\beta_0+\beta_1 r_{t-1}^2+\beta_2 r_{t-2}^2 .
\]

\textbf{(b)} By the \emph{law of total variance},
\[
\Var(r_t)=\E\bigl[\Var(r_t\mid\F_{t-1})\bigr]+\Var\bigl(\E[r_t\mid\F_{t-1}]\bigr)
=\E[\sigma_t^2]+\Var(0)=\E[\sigma_t^2].
\]
Taking the (unconditional) expectation of the variance recursion and using \emph{linearity of expectation},
\[
\E[\sigma_t^2]=\beta_0+\beta_1\E[r_{t-1}^2]+\beta_2\E[r_{t-2}^2].
\]
Because $\E[r_{t-j}]=0$ we have $\E[r_{t-j}^2]=\Var(r_{t-j})$, and by \emph{stationarity}
$\Var(r_t)=\Var(r_{t-1})=\Var(r_{t-2})=:v$. Hence
\[
v=\beta_0+\beta_1 v+\beta_2 v
\;\Longrightarrow\;
(1-\beta_1-\beta_2)\,v=\beta_0
\;\Longrightarrow\;
\boxed{\;\Var(r_t)=\frac{\beta_0}{1-\beta_1-\beta_2}\;}\quad(\beta_1+\beta_2<1).
\]
Numerically, with $\beta_0=0.2,\ \beta_1=0.3,\ \beta_2=0.5$:
\[
\Var(r_t)=\frac{0.2}{1-0.3-0.5}=\frac{0.2}{0.2}=1 .
\]

\textbf{(c)} For the variance to be positive and finite we need the denominator positive, i.e.
\[
\boxed{\;\beta_1+\beta_2<1\;}\qquad(\text{with }\beta_0>0,\ \beta_1,\beta_2\ge0).
\]
Here $\beta_1+\beta_2=0.3+0.5=0.8<1$, so a stationary variance exists (and equals $1$). This is exactly the $p=2,\,q=0$ case of the general $\mathrm{GARCH}$ stationarity rule $\sum_{j}\alpha_j+\sum_{j}\beta_j<1$: an $\mathrm{ARCH}(2)$ has $\alpha$-coefficients $\beta_1,\beta_2$ and no GARCH ($\beta_j$) feedback terms, so the rule reads $\beta_1+\beta_2<1$.

% ---------------------------------------------------------------
\solproblem{5}
\textbf{(a)} \emph{Zero mean.} $\sigma_t$ is determined by the past, hence $\F_{t-1}$-measurable, and $\eta_t$ is independent of $\F_{t-1}$ with $\E[\eta_t]=0$. By the tower property,
\[
\E[r_t]=\E\bigl[\E[r_t\mid\F_{t-1}]\bigr]
=\E\bigl[\sigma_t\,\E[\eta_t\mid\F_{t-1}]\bigr]
=\E[\sigma_t\cdot0]=0 .
\]
\emph{Uncorrelated.} Fix $\tau>0$. Since $\E[r_t]=0$, $\Cov(r_{t+\tau},r_t)=\E[r_{t+\tau}r_t]$. Both $r_t$ and $\sigma_{t+\tau}$ are $\F_{t+\tau-1}$-measurable (they depend only on information up to time $t+\tau-1$), so conditioning on $\F_{t+\tau-1}$ and pulling them out of the inner expectation,
\[
\begin{aligned}
\E[r_{t+\tau}r_t]
&=\E\bigl[\E[\sigma_{t+\tau}\eta_{t+\tau}\,r_t\mid\F_{t+\tau-1}]\bigr]\\
&=\E\bigl[r_t\,\sigma_{t+\tau}\,\E[\eta_{t+\tau}\mid\F_{t+\tau-1}]\bigr]
=\E[r_t\,\sigma_{t+\tau}\cdot0]=0 ,
\end{aligned}
\]
using $\E[\eta_{t+\tau}\mid\F_{t+\tau-1}]=\E[\eta_{t+\tau}]=0$ by independence. By symmetry of the covariance the same holds for $\tau<0$. Therefore $\Cov(r_{t+\tau},r_t)=0$ and $\Corr(r_{t+\tau},r_t)=0$ for all $\tau\neq0$. Combined with $\E[r_t]=0$ and a constant (stationary) variance, $\{r_t\}$ is \textbf{white noise}.

\textbf{(b)} As in Problem 1, $\E[r_{t-j}^2]=\Var(r_{t-j})$ (mean zero) and, using $\E[r_{t-j}^2\mid\F_{t-j-1}]=\sigma_{t-j}^2$ together with the tower property, $\E[r_{t-j}^2]=\E[\sigma_{t-j}^2]$; by stationarity all these equal $v:=\Var(r_t)$. Taking unconditional expectations of the recursion,
\[
\E[\sigma_t^2]=\alpha_0+\sum_{j=1}^{m}\alpha_j\,\E[r_{t-j}^2]+\sum_{j=1}^{r}\beta_j\,\E[\sigma_{t-j}^2]
\;\Longrightarrow\;
v=\alpha_0+\Bigl(\sum_{j}\alpha_j+\sum_{j}\beta_j\Bigr)v .
\]
Solving,
\[
\boxed{\;\Var(r_t)=v=\frac{\alpha_0}{\,1-\sum_{j}\alpha_j-\sum_{j}\beta_j\,}\;},
\]
which is positive and finite \textbf{iff}
\[
\boxed{\;\sum_{j=1}^{m}\alpha_j+\sum_{j=1}^{r}\beta_j<1\;}\qquad(\text{the weak-stationarity condition}).
\]

\textbf{(c)} For $\mathrm{GARCH}(1,1)$ with $\alpha_0=2\times10^{-5},\ \alpha_1=0.10,\ \beta_1=0.85$:
\[
\alpha_1+\beta_1=0.10+0.85=0.95<1,
\]
so the process is weakly stationary. Then
\[
\Var(r_t)=\frac{\alpha_0}{1-\alpha_1-\beta_1}
=\frac{2\times10^{-5}}{1-0.95}
=\frac{2\times10^{-5}}{0.05}
=4\times10^{-4},
\]
and the implied daily volatility is
\[
\sqrt{\Var(r_t)}=\sqrt{4\times10^{-4}}=2\times10^{-2}=2\%\ \text{per day},
\]
a realistic equity figure. The persistence $\alpha_1+\beta_1=0.95$ is close to $1$, the hallmark of financial GARCH fits: volatility shocks decay only slowly (sticky volatility, long quiet/turbulent regimes), yet the process stays stationary because the sum is still strictly below $1$.

% ---------------------------------------------------------------
\solproblem{5}
\textbf{(a)} Fix $\tau>0$ and start from the covariance, using stationarity ($\E[r_{t+\tau}^2]=\E[r_t^2]$):
\[
\Cov(r_{t+\tau}^2,r_t^2)=\E[r_{t+\tau}^2 r_t^2]-\bigl(\E[r_t^2]\bigr)^2 .
\]
Condition on $\F_{t+\tau-1}$: both $r_t^2$ and $\sigma_{t+\tau}^2=\alpha_0+\alpha_1 r_{t+\tau-1}^2$ are measurable there, and $\E[\eta_{t+\tau}^2\mid\F_{t+\tau-1}]=\E[\eta_{t+\tau}^2]=1$, so
\[
\E[r_{t+\tau}^2 r_t^2\mid\F_{t+\tau-1}]
=r_t^2\,\sigma_{t+\tau}^2\,\E[\eta_{t+\tau}^2\mid\F_{t+\tau-1}]
=r_t^2\bigl(\alpha_0+\alpha_1 r_{t+\tau-1}^2\bigr).
\]
Taking expectations,
\[
\E[r_{t+\tau}^2 r_t^2]=\alpha_0\,\E[r_t^2]+\alpha_1\,\E[r_{t+\tau-1}^2 r_t^2].
\]
Write $C_\tau:=\Cov(r_{t+\tau}^2,r_t^2)$ and use $\E[r_{t+\tau-1}^2 r_t^2]=C_{\tau-1}+\bigl(\E[r_t^2]\bigr)^2$ (again by stationarity, $\E[r_{t+\tau-1}^2]=\E[r_t^2]$). Then
\[
C_\tau+\bigl(\E[r_t^2]\bigr)^2
=\alpha_0\,\E[r_t^2]+\alpha_1\Bigl(C_{\tau-1}+\bigl(\E[r_t^2]\bigr)^2\Bigr),
\]
i.e.
\[
C_\tau=\alpha_1 C_{\tau-1}
+\underbrace{\alpha_0\,\E[r_t^2]+\alpha_1\bigl(\E[r_t^2]\bigr)^2-\bigl(\E[r_t^2]\bigr)^2}_{=:R}.
\]
Substitute $\E[r_t^2]=\Var(r_t)=\alpha_0/(1-\alpha_1)$:
\[
R=\alpha_0\cdot\frac{\alpha_0}{1-\alpha_1}-(1-\alpha_1)\Bigl(\frac{\alpha_0}{1-\alpha_1}\Bigr)^2
=\frac{\alpha_0^2}{1-\alpha_1}-\frac{\alpha_0^2}{1-\alpha_1}=0 .
\]
Hence
\[
\boxed{\;C_\tau=\alpha_1\,C_{\tau-1}\;}\qquad(\tau>0).
\]
Iterating $\tau$ times from $C_0=\Var(r_t^2)$ gives $C_\tau=\alpha_1^{\tau}\,\Var(r_t^2)$, so
\[
\Corr(r_{t+\tau}^2,r_t^2)=\frac{C_\tau}{\Var(r_t^2)}=\alpha_1^{\tau},
\]
and by symmetry of the autocovariance, $\Corr(r_{t+\tau}^2,r_t^2)=\alpha_1^{\abs\tau}$ for all $\tau$. (The condition $0<\alpha_1<1/\sqrt3$, i.e.\ $\alpha_1^2<1/3$, guarantees $\Var(r_t^2)=\E[r_t^4]-(\E[r_t^2])^2<\infty$, so the correlation is well defined.)

\textbf{(b)} With $\alpha_1=0.5$, $\alpha_1^2=0.25<1/3$, so the fourth moment exists and
\[
\kappa=\frac{3(1-\alpha_1^2)}{1-3\alpha_1^2}
=\frac{3(1-0.25)}{1-0.75}
=\frac{3\cdot0.75}{0.25}
=\frac{2.25}{0.25}=9 .
\]
Since $\kappa=9>3$, the marginal law has \emph{heavier tails than the normal} (excess kurtosis $9-3=6$): ARCH generates fat tails endogenously, even with Gaussian innovations. The squared-return autocorrelations are $\Corr(r_{t+\tau}^2,r_t^2)=0.5^{\abs\tau}$:
\[
\tau=1:\ 0.5,\qquad \tau=2:\ 0.25,\qquad \tau=3:\ 0.125 .
\]
The two thresholds differ because they control different moments: $\alpha_1<1$ keeps the \emph{second} moment ($\Var r_t=\alpha_0/(1-\alpha_1)$) finite, whereas the \emph{fourth} moment involves $\E[\sigma_t^4]$ and the recursion for it carries the factor $3\alpha_1^2$ (from $\E[\eta^4]=3$); a finite solution needs $3\alpha_1^2<1$, i.e.\ the strictly tighter $\alpha_1^2<1/3$. As $\alpha_1^2\uparrow1/3$ the kurtosis $\to\infty$.

\textbf{(c)} With $\alpha_0=4\times10^{-4}$, $\alpha_1=0.5$ and $r_t=0.04$ (so $r_t^2=1.6\times10^{-3}$),
\[
\sigma_{t+1}^2=\alpha_0+\alpha_1 r_t^2
=4\times10^{-4}+0.5\,(1.6\times10^{-3})
=4\times10^{-4}+8\times10^{-4}=1.2\times10^{-3},
\]
so $\sigma_{t+1}=\sqrt{1.2\times10^{-3}}\approx0.0346$. Since $r_{t+1}\mid\F_t\sim\mathcal N(0,\sigma_{t+1}^2)$, the $95\%$ conditional predictive interval is
\[
0\pm1.96\,\sigma_{t+1}=\pm1.96\,(0.0346)\approx\pm0.0679 .
\]
The unconditional variance is $\Var(r_t)=\alpha_0/(1-\alpha_1)=4\times10^{-4}/0.5=8\times10^{-4}$, so $\sqrt{\Var(r_t)}\approx0.0283$ and the unconditional band is $\pm1.96\,(0.0283)\approx\pm0.0554$. The conditional band ($\pm0.068$) is markedly \emph{wider} than the unconditional one ($\pm0.055$): after the large move $r_t=0.04$ the model raises its volatility forecast --- exactly the volatility-clustering behaviour ARCH is designed to capture.

% ---------------------------------------------------------------
\solproblem{5}
\textbf{(a)} With $n=100$, $n(n+2)=100\cdot102=10200$, and $m=5$:
\[
Q_5=10200\left(\frac{0.05^2}{99}+\frac{(-0.12)^2}{98}+\frac{0.08^2}{97}+\frac{(-0.10)^2}{96}+\frac{0.06^2}{95}\right).
\]
Term by term,
\[
\begin{aligned}
\frac{0.0025}{99}&=2.525\times10^{-5}, &
\frac{0.0144}{98}&=1.469\times10^{-4}, &
\frac{0.0064}{97}&=6.598\times10^{-5},\\
\frac{0.0100}{96}&=1.042\times10^{-4}, &
\frac{0.0036}{95}&=3.789\times10^{-5}. &&
\end{aligned}
\]
Summing the five terms gives $3.802\times10^{-4}$, so
\[
Q_5\approx10200\times3.802\times10^{-4}\approx \boxed{3.88}.
\]

\textbf{(b)} For an $\mathrm{ARMA}(1,1)$, $p=q=1$, so $\mathrm{df}=m-p-q=5-1-1=3$. The $5\%$ critical value is $\chi^2_{3,0.95}=7.815$. Since $Q_5\approx3.88<7.815$, we \textbf{do not reject} $H_0$: there is no significant evidence of residual autocorrelation up to lag $5$, so the $\mathrm{ARMA}(1,1)$ appears \emph{adequate}.

Had an $\mathrm{AR}(1)$ been fitted, $p=1,q=0$, so $\mathrm{df}=5-1-0=4$ and the critical value rises to $\chi^2_{4,0.95}=9.488$. Since $3.88<9.488$, the decision is unchanged (do not reject). In general fewer fitted parameters leave more degrees of freedom and widen the acceptance region.

\textbf{(c)} The parameter counts are $k=p+q+1$: for $\mathrm{ARMA}(1,1)$, $k=3$; for $\mathrm{AR}(3)$, $k=4$ (one extra parameter). With $n=100$, $\log n=\log 100\approx4.605$. Comparing penalties when adding the one extra parameter:
\[
\begin{aligned}
\Delta(\text{penalty})_{\mathrm{AIC}}&=2\,(k_{\text{big}}-k_{\text{small}})=2,\\
\Delta(\text{penalty})_{\mathrm{BIC}}&=(\log n)(k_{\text{big}}-k_{\text{small}})=\log 100\approx4.605 .
\end{aligned}
\]
So the $\mathrm{AR}(3)$ is preferred only if $-2\log L$ \emph{drops} by more than the extra penalty:
\[
\text{AIC: drop }>2\ \ (\text{i.e.}\ \log L\ \text{improves by}>1);\qquad
\text{BIC: drop }>4.605\ \ (\log L\ \text{improves by}>2.30).
\]
Since BIC's $\log n=4.605>2$ charges more per parameter (as $n>e^2\approx7.4$), \textbf{AIC is more willing to select the larger $\mathrm{AR}(3)$ model}; BIC is harder to satisfy and favours the parsimonious $\mathrm{ARMA}(1,1)$.

In-sample $R^2=1-s_e^2/s_y^2$ is a poor selection tool because adding parameters can only \emph{increase} it (the larger nested model fits the training data at least as well), so a higher $R^2$ for $\mathrm{AR}(3)$ is no evidence of a genuinely better model; moreover even the true model has a bounded $R^2$ (a stochastic process is not perfectly predictable). Information criteria penalise complexity and are the correct basis for the comparison.

\end{document}
